\section*{Chapitre \ref{ch:b1:radiometric-methods} --- Les méthodes de datation radiométrique}
% values: calc/ch25.py

\begin{solution}{exo:b1:radiometric-methods:1}
(a)~Ici $D^*/P = 2^n - 1 = 3$ donne $n = 2$~: deux demi-vies. (b)~Rubidium~87~:
$e^{\lambda t} - 1 = e^{\num{1.3972e-11} \times \num{3.10e8}} - 1 = 0.004341$~;
uranium~235~: $e^{\num{9.8485e-10} \times \num{3.10e8}} - 1 = 0.3570$. Le
rapport de l'uranium est environ 80 fois plus grand~: une erreur absolue de mesure donnée en
est une part bien plus petite. (c)~$2^{-n} < 0.01$ quand $n > \log_2 100 = 6.64$~: après
environ sept demi-vies.
\end{solution}

\begin{solution}{exo:b1:radiometric-methods:2}
(a)~Radiocarbone sur le charbon de bois. (b)~Uranium--plomb sur le zircon, ou
argon--argon (potassium--argon) sur la sanidine, \omterm{def:b1:rock-forming-minerals:mineral}{minéral} riche en potassium.
(c)~Uranium--plomb sur zircon, qui survit à une longue histoire~; rubidium--strontium
sur roches totales est possible~; potassium--argon risquerait de donner des âges abaissés par l'argon
perdu lors d'un chauffage ultérieur. (d)~Radiocarbone sur le carbonate de la coquille, près de la
limite de la méthode (environ \num{70000} ans). (e)~Aucune directement~: son âge est
encadré par des roches volcaniques datées au-dessus et en dessous, ou trouvé en corrélant ses
fossiles avec une coupe datée (\cref{ch:b1:stratigraphic-scale}).
\end{solution}

\begin{solution}{exo:b1:radiometric-methods:3}
(a)~$t = -8033 \ln 0.25 = 8033 \times 1.3863 = \num{11136}$ ans BP.
(b)~La vie moyenne est $5700/\ln 2 = \num{8223}$ ans~: $t = 8223 \times 1.3863 =
\num{11400}$ ans. (c)~Les deux supposent que la proportion de carbone~14 dans le carbone vivant
était alors la même que dans l'étalon~; ce n'était pas le cas. L'âge calendaire se lit
sur une \omterm{def:b1:radiometric-methods:radiocarbon}{courbe d'étalonnage} telle qu'IntCal20, comme un intervalle d'années cal BP.
\end{solution}

\begin{solution}{exo:b1:radiometric-methods:4}
(a)~Chaque horloge démarre quand son \omterm{def:b1:rock-forming-minerals:mineral}{minéral} se refroidit sous sa \omterm{def:b1:radiometric-methods:closure}{température de fermeture}~:
le zircon (au-dessus de \qty{750}{\celsius}) se ferme en premier, la hornblende (environ
\qty{530}{\celsius}) ensuite, la biotite (environ \qty{280}{\celsius}) en dernier, si bien que les âges
décroissent dans cet ordre à mesure que le granite se refroidissait. (b)~$(530 - 280)/(410 - 398) =
250/12 = \qty{21}{\celsius}$ par million d'années. (c)~Le zircon s'est fermé au-dessus de
\qty{750}{\celsius}, donc le granite s'est refroidi d'au moins $750 - 530 =
\qty{220}{\celsius}$ en \qty{10}{\Ma}~: au moins \qty{22}{\celsius} par million
d'années, une borne inférieure.
\end{solution}

\begin{solution}{exo:b1:radiometric-methods:5}
(a)~$\bar x = 7.5$, $\bar y = 0.809525$, $S_{xx} = 81.0$, $S_{xy} = 1.13965$~;
$m = 1.13965/81.0 = 0.0140698$, $b = 0.809525 - 7.5 \times 0.0140698 = 0.70400$.
(b)~$t = \ln(1.0140698)/\num{1.3972e-11} = \qty{1000.0}{\Ma}$~; \omterm{def:b1:radiometric-methods:isochron}{rapport initial}
$0.7040$. (c)~Les résidus sont $+1.6$, $-2.5$, $+0.7$ et $+0.2$ (en unités de
\num{e-4})~: $s = \sqrt{\num{9.34e-8}/2} = 0.00022$, $\sigma_m = 0.00022/\sqrt{81}
= \num{2.40e-5}$ et $\sigma_t = \num{2.40e-5}/(\num{1.3972e-11} \times 1.0141) =
\qty{1.7}{\Ma}$~: \qty[uncertainty-mode=separate]{1000.0(17)}{\Ma}. (d)~$0.7040$
est proche de la valeur du \omterm{def:b1:layered-earth:layers}{manteau} $0.7025$ et bien en dessous de celle de la \omterm{def:b1:layered-earth:layers}{croûte} $0.710$~: une source
surtout dans le \omterm{def:b1:layered-earth:layers}{manteau}.
\end{solution}

\begin{solution}{exo:b1:radiometric-methods:6}
(a)~$0.100/39.098 \times \num{1.17e-4} = \qty{2.992e-7}{mol}$ de potassium~40 par
gramme. (b)~${}^{40}\mathrm{Ar}^*/{}^{40}\mathrm{K} = \num{1.169e-9}/\num{2.992e-7}
= 0.0039065$~; $\lambda/\lambda_e = 5.543/0.581 = 9.5404$~; $t = \ln(1 + 9.5404
\times 0.0039065)/\num{5.543e-10} = \ln(1.037269)/\num{5.543e-10} =
\qty{66.01}{\Ma}$. (c)~L'argon radiogénique serait $0.995 \times 1.169$
nanomoles, ce qui donne \qty{65.69}{\Ma}~: l'âge mesuré serait trop vieux de \qty{0.32}{\Ma},
dix fois les \num{32000} ans dans lesquels l'impact et
l'extinction de la fin du Crétacé se révèlent désormais coïncider. L'argon en excès est une préoccupation surtout dans les roches jeunes et les roches totales~; un
\omterm{def:b1:rock-forming-minerals:mineral}{minéral} tel que la sanidine ne doit pas prendre d'argon en cristallisant, et un
spectre d'âge argon--argon le révélerait.
\end{solution}

\begin{solution}{exo:b1:radiometric-methods:7}
(a)~Environ \num{2000} cal BP~: la courbe passe par \num{2070} ans BP à cet endroit.
(b)~D'environ \num{2495} à \num{2699} cal BP, deux cents ans~: la courbe est
presque plate entre \num{2500} et \num{2700} cal BP, si bien que chaque âge calendaire
de cet intervalle a presque le même \omterm{def:b1:radiometric-methods:radiocarbon}{âge radiocarbone}, et qu'une mesure plus précise
ne réduirait pas l'intervalle. (c)~À \num{2000} cal BP l'\omterm{def:b1:radiometric-methods:radiocarbon}{âge radiocarbone}
est de \num{70} ans \emph{plus vieux} que l'âge calendaire~; à \num{8000} cal BP
\num{822} ans plus jeune (\qty{10}{\percent})~; à \num{12000} cal BP \num{1694}
ans plus jeune (\qty{14}{\percent}). La différence n'est pas proportionnelle à
l'âge~: elle change même de signe. Elle ne peut venir de la seule demi-vie, qui
donnerait un \qty{2.4}{\percent} constant.
\end{solution}

\begin{solution}{exo:b1:radiometric-methods:8}
(a)~Le banc de cendre en dessous (\qty{94.50}{\Ma}) donne un âge maximal, la coulée de lave au-dessus
(\qty{93.80}{\Ma}) un âge minimal~: la limite est encadrée entre eux. Le dyke est
aussi un minimum, mais plus faible, plus jeune que la lave. Le zircon détritique est un
maximum seulement pour le grès, et plus faible que le banc de cendre, qui se trouve en dessous du
grès. (b)~$f = (12 - 4)/(20 - 4) = 0.5$~; $t = 94.50 - 0.5 \times 0.70 =
\qty{94.15}{\Ma}$~; $\sigma_t = \sqrt{0.25 \times 0.10^2 + 0.25 \times 0.12^2} =
\qty{0.08}{\Ma}$~: \qty[uncertainty-mode=separate]{94.15(8)}{\Ma}. (c)~Ses
grains viennent d'une roche d'au moins \qty{95.3}{\Ma}, érodée quelque part dans la
zone source~; cela montre que le grès est plus jeune que \qty{95.3}{\Ma}, ce que
le banc de cendre en dessous prouve déjà plus étroitement.
\end{solution}

\begin{solution}{exo:b1:radiometric-methods:9}
(a)~Grain~A~: $\ln(1.08401)/\num{1.55125e-10} = \qty{520.0}{\Ma}$ et
$\ln(1.6688)/\num{9.8485e-10} = \qty{520.0}{\Ma}$. Grain~B~: \qty{368.4}{\Ma} et
\qty{389.9}{\Ma}. (b)~Le grain~A est concordant. (c)~$0.05881/0.08401 = 0.7000$ et
$0.4682/0.6688 = 0.7001$. L'âge est $\ln(1 + R)/\lambda$, non proportionnel à
$R$~: pour le petit rapport de l'uranium~238 $\ln(1+R) \approx R$ et l'âge chute de
près de \qty{30}{\percent} (à \qty{368}{\Ma})~; pour le plus grand rapport de l'uranium~235
le logarithme chute moins, et l'âge seulement à \qty{390}{\Ma}. Les deux sont trop
jeunes, et ils ne s'accordent pas~: le grain est discordant.
\end{solution}

\begin{solution}{exo:b1:radiometric-methods:10}
(a)~Avec $u = \lambda t$, $D^*/P = e^u - 1 = u + u^2/2 + \dots$, donc $(D^*/P)/\lambda
= t\,(1 + u/2 + \dots)$~: l'approximation surestime l'âge d'une erreur relative
voisine de $u/2$. (b)~Rubidium~87 à \qty{310}{\Ma}~: $u = 0.00433$, erreur
\qty{0.22}{\percent} (\qty{0.7}{\Ma})~; potassium~40 à \qty{296}{\Ma}~: $u =
0.164$, erreur \qty{8.7}{\percent} (\qty{26}{\Ma}~; $u/2$ donne \qty{8.2}{\percent})~;
uranium~238 à \qty{4.4}{\Ga}~: $u = 0.683$, erreur \qty{43}{\percent} ---
l'approximation et même son estimation d'erreur ($u/2 = \qty{34}{\percent}$) échouent.
(c)~$R = 0.004341$, $\sigma_R = \num{4.34e-6}$~; d'après la
\cref{prop:b1:radiometric-methods:precision}, $\sigma_t =
\num{4.34e-6}/(\num{1.3972e-11} \times 1.0043) = \qty{0.31}{\Ma}$.
\end{solution}

\begin{solution}{exo:b1:radiometric-methods:11}
(a)~Soit ${}^{40}\mathrm{K}$ le potassium~40 d'aujourd'hui. Au temps $t_2$ avant
le présent il y avait ${}^{40}\mathrm{K}\,e^{\lambda t_2}$ atomes, et l'argon
fait entre $t_1$ et $t_2$ était, d'après la \cref{prop:b1:radiometric-methods:kar} appliquée
à cet intervalle, $(\lambda_e/\lambda)\,{}^{40}\mathrm{K}\,e^{\lambda
t_2}\bigl(e^{\lambda (t_1 - t_2)} - 1\bigr) = (\lambda_e/\lambda)\,
{}^{40}\mathrm{K}\,\bigl(e^{\lambda t_1} - e^{\lambda t_2}\bigr)$. Une fraction
$1 - f$ en est restée. Depuis $t_2$, le \omterm{def:b1:rock-forming-minerals:mineral}{minéral} fermé a fait
$(\lambda_e/\lambda)\,{}^{40}\mathrm{K}\,(e^{\lambda t_2} - 1)$ de plus. En ajoutant et en
multipliant par $\lambda/(\lambda_e\,{}^{40}\mathrm{K})$ on obtient l'expression.
(b)~L'âge apparent est $\ln(1 + \text{expression})/\lambda$~: \qty{400.0}{\Ma} pour
$f = 0$, \qty{342.3}{\Ma} pour $0.3$, \qty{282.7}{\Ma} pour $0.6$ et
\qty{200.0}{\Ma} pour $1$. (c)~Tout âge compris entre \qty{200}{\Ma} et
\qty{400}{\Ma} peut être produit par une perte partielle, et un nombre unique ne peut
dire s'il enregistre un événement réel à cette date ou un mélange~; un spectre d'âge,
ou d'autres horloges dans la même roche, sont nécessaires.
\end{solution}

\begin{solution}{exo:b1:radiometric-methods:12}
(a)~L'âge est proportionnel à la vie moyenne employée~: $10306 \times 8223/8033 =
\num{10550}$ ans. (b)~$10550 - 10306 = 244$ ans, environ un septième des
\num{1694}. (c)~Si l'activité initiale était $k A_0$, le vrai âge est $t = \tau
\ln(k A_0/A) = \num{10550} + \tau \ln k$~; avec $t = \num{12000}$, $\ln k =
1450/8223 = 0.176$ et $k = 1.19$~: environ \qty{19}{\percent} de carbone~14 de plus que
dans l'étalon.
\end{solution}

\begin{solution}{pb:b1:radiometric-methods:1}
\begin{enumerate}
  \item Les points montent régulièrement avec $x$~; leurs écarts à une droite sont au plus
        de deux unités à la quatrième décimale, pour une étendue de $0.04$ en $y$~: ils sont sur une
        droite à faible dispersion près.
  \item $\bar x = 4.90333$, $\bar y = 0.729288$, $S_{xx} = 62.1253$, $S_{xy} =
        0.269335$~; $m = 0.269335/62.1253 = 0.0043354$.
  \item $b = 0.729288 - 0.0043354 \times 4.90333 = 0.70803$~: le
        $^{87}\mathrm{Sr}/^{86}\mathrm{Sr}$ du granite quand il a cristallisé,
        partagé par toutes ses parties.
  \item $t = \ln(1.0043354)/\num{1.3972e-11} = \qty{309.6}{\Ma}$.
  \item Les résidus sont $+9.1$, $-17.2$, $+6.9$, $-6.4$, $+15.1$ et $-7.5$ (en
        unités de \num{e-5})~; $s = 0.000137$, $\sigma_m = 0.000137/\sqrt{62.1253} =
        \num{1.74e-5}$, $\sigma_t = \num{1.74e-5}/(\num{1.3972e-11} \times 1.0043)
        = \qty{1.2}{\Ma}$~: \qty[uncertainty-mode=separate]{309.6(12)}{\Ma}, et
        $\sigma_b = 0.00010$.
  \item $0.7080$ est plus proche de la valeur crustale $0.710$ que de celle du \omterm{def:b1:layered-earth:layers}{manteau}, $0.7025$.
        Pour un mélange, $(0.70803 - 0.7025)/(0.710 - 0.7025) = 0.74$~: environ trois
        quarts du strontium venaient de \omterm{def:b1:layered-earth:layers}{croûte} fondue.
  \item $0.0740/39.098 \times \num{1.17e-4} = \qty{2.214e-7}{mol}$ par gramme,
        multiplié par $\num{6.022e23}$~: \num{1.33e17} atomes.
  \item $^{40}\mathrm{Ar}^*/^{40}\mathrm{K} = \num{4.139e-9}/\num{2.2144e-7} =
        0.018691$~; $1 + 9.5404 \times 0.018691 = 1.17832$~; $t =
        \ln(1.17832)/\num{5.543e-10} = \qty{296.0}{\Ma}$.
  \item $^{40}\mathrm{K} = \qty{3.292e-8}{mol.g^{-1}}$, rapport $0.019239$,
        $\ln(1.18355)/\num{5.543e-10} = \qty{304.0}{\Ma}$.
  \item $\ln(1.018691)/\num{5.543e-10} = \qty{33.4}{\Ma}$, neuf fois trop jeune,
        parce qu'une désintégration sur dix seulement donne de l'argon. Le calcium~40 est
        l'isotope courant du calcium, abondant dans les roches~: le peu qui est fait par désintégration ne peut
        être distingué du calcium que les minéraux contenaient dès le départ.
  \item L'argon en excès décale les âges de l'ordre d'un million d'années, ce qui
        compte pour les roches jeunes~; ici il serait bien en dessous de \qty{1}{\percent} de
        l'âge. Et les cristaux de \omterm{def:b1:rock-forming-minerals:groups}{mica} et d'\omterm{def:b1:rock-forming-minerals:groups}{amphibole} ne prennent presque pas d'argon quand ils
        se forment.
  \item Soit le granite s'est refroidi lentement et les \omterm{def:b1:rock-forming-minerals:groups}{micas} et l'\omterm{def:b1:rock-forming-minerals:groups}{amphibole} se sont fermés après
        la cristallisation (\omterm{def:b1:radiometric-methods:closure}{âges de refroidissement}), soit un chauffage ultérieur a chassé une partie de leur
        argon. Un plateau plat dans un spectre d'âge argon--argon soutient un \omterm{def:b1:radiometric-methods:closure}{âge de
        refroidissement}~; un spectre sans plateau avertit que les minéraux ont perdu de l'argon.
        L'ordre des âges, la hornblende avant la biotite comme le prédisent leurs \omterm{def:b1:radiometric-methods:closure}{températures de
        fermeture}, indique aussi un refroidissement.
  \item $\ln(1.049361)/\num{1.55125e-10} = \qty{310.6}{\Ma}$ et
        $\ln(1.35784)/\num{9.8485e-10} = \qty{310.6}{\Ma}$~: concordants.
  \item Zircon \qty{310.6}{\Ma}, fermeture au-dessus de \qty{750}{\celsius}~:
        cristallisation~; l'\omterm{def:b1:radiometric-methods:isochron}{isochrone} sur roche totale,
        \qty[uncertainty-mode=separate]{309.6(12)}{\Ma}, la date aussi,
        et s'accorde dans son incertitude~; hornblende \qty{304.0}{\Ma}~: refroidissement
        à travers environ \qty{530}{\celsius}~; biotite \qty{296.0}{\Ma}~: refroidissement
        à travers environ \qty{280}{\celsius}. Le graphe descend de plus de
        \qty{750}{\celsius} à \qty{310.6}{\Ma} à \qty{280}{\celsius} à
        \qty{296.0}{\Ma}.
  \item $(530 - 280)/(304.02 - 296.03) = 250/7.99 = \qty{31}{\celsius}$ par
        million d'années.
  \item $(750 - 530)/(310.60 - 304.02) = 220/6.58 = \qty{33}{\celsius}$ par
        million d'années au moins, puisque le zircon s'est fermé au-dessus de
        \qty{750}{\celsius}~: le refroidissement était au moins aussi rapide au début.
  \item La différence de températures, $250 \pm \sqrt{40^2 + 40^2} = 250 \pm
        57$, donne \qtyrange{24}{38}{\celsius} par million d'années.
  \item $(550 - 350)/7.99 = \qty{25}{\celsius}$ par million d'années, dans l'intervalle
        de la question~17~: l'incertitude vient des \omterm{def:b1:radiometric-methods:closure}{températures de fermeture}, non
        des âges.
  \item $(280 - 10)/30 = \qty{9.0}{km}$.
  \item $\ln(1.039108)/\num{1.55125e-10} = \qty{247.30}{\Ma}$, et de même d'après
        le rapport de l'uranium~235~: concordants.
  \item Du granite, érodé à la \omterm{def:b1:cross-sections:unconformity}{discordance}~: leur âge est son
        âge de cristallisation. Ils montrent seulement que le grès est plus jeune que
        \qty{310}{\Ma}~: un âge maximal, bien plus faible que celui du banc de cendre.
  \item $(58 - 18)/(247.30 - 245.30) = \qty{20.0}{m}$ par million d'années~; la base
        est plus vieille que le banc de cendre inférieur de $18/20.0 = \qty{0.90}{\Ma}$~:
        \qty{248.20}{\Ma}.
  \item De \qty{309.6}{\Ma} (\omterm{def:b1:radiometric-methods:isochron}{isochrone}) ou \qty{310.6}{\Ma} (zircon) à
        \qty{248.2}{\Ma}~: environ \qty{62}{\Ma}.
  \item \qty{9.0}{km} en $296.0 - 248.2 = \qty{47.8}{\Ma}$~: \qty{0.19}{km} par
        million d'années, \qty{0.19}{mm} par an. Cela suppose que la biotite
        datée venait de juste sous la \omterm{def:b1:cross-sections:unconformity}{discordance}, si bien que toute la roche au-dessus
        d'elle dans le granite et le schiste avait été enlevée quand le sable est arrivé,
        et que le \omterm{def:b1:heat-budget:geotherm}{géotherme} est resté constant.
  \item Le granite s'est refroidi d'environ \qty{31}{\celsius} par million d'années
        entre les fermetures de la hornblende et de la biotite.
\end{enumerate}
\end{solution}
